\subsection{Longitudinal energy bands}
\begin{figure}[H]\centering
\begin{tikzpicture}
\begin{axis}[width=.92\linewidth,height=6.8cm,
xlabel={Source-to-drain position, $X$ ($\mu$m)},
ylabel={$E_C$ (eV)},xmin=-.85,xmax=.85,
ymin=-1.5,ymax=.55,grid=both,legend pos=south west]
\addplot+[no marks,thick]
table[col sep=comma,x=X_um,y=Ec_eV]
{data/part5_longitudinal_tsi30.csv};
\addlegendentry{$T_{\mathrm{SI}}=30$ nm}
\addplot+[no marks,thick,dashed]
table[col sep=comma,x=X_um,y=Ec_eV]
{data/part5_longitudinal_tsi10.csv};
\addlegendentry{$T_{\mathrm{SI}}=10$ nm}
\end{axis}
\end{tikzpicture}
\caption{Conduction-band energy, sampled 1.5 nm beneath the
gate-oxide interface at $V_{GS}=0$, $V_{DS}=1.2$ V.
The source lies on the left and the drain on the right.
Shown are genuine simulator points selected for plotting.}
\label{fig:part5_long_ec}
\end{figure}

\begin{figure}[H]\centering
\begin{tikzpicture}
\begin{axis}[width=.92\linewidth,height=6.8cm,
xlabel={Source-to-drain position, $X$ ($\mu$m)},
ylabel={$E_V$ (eV)},xmin=-.85,xmax=.85,
ymin=-2.35,ymax=-.55,grid=both,legend pos=south west]
\addplot+[no marks,thick]
table[col sep=comma,x=X_um,y=Ev_eV]
{data/part5_longitudinal_tsi30.csv};
\addlegendentry{$T_{\mathrm{SI}}=30$ nm}
\addplot+[no marks,thick,dashed]
table[col sep=comma,x=X_um,y=Ev_eV]
{data/part5_longitudinal_tsi10.csv};
\addlegendentry{$T_{\mathrm{SI}}=10$ nm}
\end{axis}
\end{tikzpicture}
\caption{Valence-band energy on the same two longitudinal
cutlines and at identical terminal biases.}
\label{fig:part5_long_ev}
\end{figure}

\paragraph{Source-channel barrier.}
With the source conduction band referenced at
$X\simeq-0.8\,\mu$m, $E_C$ equals $-0.16156$ eV
(30 nm) or $-0.15609$ eV (10 nm).
The maximum $E_C$ along each line occurs at
$X=-0.1275\,\mu$m and $-0.1343\,\mu$m,
respectively, with values $0.39918$ and
$0.38485$ eV. Hence the extracted local
source-to-crest barrier heights are
$0.56074$ and $0.54094$ eV. In particular,
the thin-film cutline barrier is about
20 meV \emph{lower} than the thick-film
one under these specific bias conditions.
It would be incorrect to claim a taller
thin-film band-edge crest from these results,
even though the measured thin-film
$I_{\mathrm{OFF}}$ is lower (Part~2).
Leakage depends on the complete two-dimensional
transport path rather than one horizontal
surface-adjacent line.

\paragraph{Drain-side slope and electrostatic lowering.}
From $X\simeq0.20$ to $0.40\,\mu$m,
$E_C$ falls by 13.9 meV for 30 nm
and by 7.2 meV for 10 nm, before
a steeper descent through the drain junction.
This particular pre-junction band-edge slope
is smaller for the thinner film.
The source/drain $E_C$ plateaus differ by
approximately 1.20 eV, consistent with
$V_{DS}=1.2$ V. Drain electrostatic coupling
can lower the source-channel barrier in short
devices, but \emph{quantitative drain-induced
barrier lowering} requires a corresponding
low-$V_{DS}$ band snapshot at the same gate
bias. The present two TDRs both use
$V_{DS}=1.2$ V, so this study does not
invent a numerical change in barrier height
between drain biases.
